\documentclass[10pt,twocolumn,letterpaper]{article}

% --- Geometry ---
\usepackage[
  top=0.75in, bottom=1in, left=0.75in, right=0.75in,
  columnsep=0.3in
]{geometry}

% --- Fonts and encoding ---
\usepackage[T1]{fontenc}
\usepackage{mathptmx}
\usepackage{courier}
\usepackage[scaled=0.92]{helvet}

% --- Math and symbols ---
\usepackage{amsmath,amssymb,amsfonts}

% --- Tables and figures ---
\usepackage{booktabs}
\usepackage{tabularx}
\usepackage{array}

% --- Bibliography ---
\usepackage[numbers,sort&compress]{natbib}

% --- Hyperlinks ---
\usepackage[
  colorlinks=true,
  linkcolor=black,
  citecolor=black,
  urlcolor=blue!70!black,
  bookmarks=false
]{hyperref}

% --- Misc ---
\usepackage{enumitem}
\usepackage{xcolor}
\usepackage{microtype}
\usepackage{balance}

% --- Section formatting ---
\usepackage{titlesec}
\titleformat{\section}{\normalfont\large\bfseries}{\thesection.}{0.5em}{}
\titleformat{\subsection}{\normalfont\normalsize\bfseries}{\thesubsection}{0.5em}{}
\titlespacing*{\section}{0pt}{10pt plus 2pt minus 2pt}{4pt plus 1pt}
\titlespacing*{\subsection}{0pt}{8pt plus 2pt minus 1pt}{2pt plus 1pt}

% --- Title formatting ---
\makeatletter
\renewcommand{\maketitle}{%
  \twocolumn[%
    \begin{center}
      {\LARGE\bfseries\@title\par}
      \vskip 1em
      {\large\@author\par}
      \vskip 0.5em
      {\normalsize\itshape\@date\par}
      \vskip 1.5em
    \end{center}
  ]%
}
\makeatother

% --- Abstract formatting ---
\renewenvironment{abstract}{%
  \small
  \begin{center}
    \textbf{Abstract}\vspace{-0.5em}
  \end{center}
  \begin{quotation}\noindent
}{%
  \end{quotation}\vspace{0.5em}\par
}

% --- Compact lists ---
\setlist{nosep,leftmargin=1.2em}

% --- Code formatting ---
\newcommand{\code}[1]{\texttt{\small #1}}

% --- Category theory shortcuts ---
\newcommand{\Sign}{\mathbf{Sign}}
\newcommand{\Sen}{\mathrm{Sen}}
\newcommand{\Mod}{\mathrm{Mod}}
\newcommand{\Cat}{\mathbf{Cat}}
\newcommand{\Set}{\mathbf{Set}}
\newcommand{\Ins}{\mathbf{Ins}}

\begin{document}

\title{Eigenius as a Grothendieck Institution:\\Structured Fibers for AI-Driven Engineering\\and Scientific Research}

\author{%
  Hans-Martin Will\\
  \texttt{hwill@acm.org}
}

\date{DRAFT --- \today}

\maketitle
\thispagestyle{empty}

% ============================================================
\begin{abstract}
Knowledge graphs and triple stores model domain data as flat elements over a schema---the category of elements of a $\Set$-valued functor.
This is adequate for database records but loses the internal structure of results produced by domain-specific reasoning systems: finite element displacement fields have mesh convergence behavior, proof terms have reduction semantics, molecular docking poses have scoring landscapes.
These are morphisms within fibers, invisible to flat storage.
We present Eigenius, an open-source platform whose capability protocol implements a Grothendieck institution: multiple logical systems (institutions in the sense of Goguen and Burstall) share a typed knowledge graph, with each institution contributing structured fibers---categories of results with their own internal morphisms---rather than flat sets of data points.
Four epistemic categories (declared, observed, derived, verified) arise from fiber membership in the total category.
We develop two worked examples: (1)~mechanical engineering, where CAD constraint solving, AI-driven generative design, and finite element stress simulation operate as three institutions over a shared model; and (2)~biopharmaceutical R\&D, where molecular docking, ADMET prediction, experimental assay data, and pharmacokinetic modeling operate as institutions over a shared compound knowledge graph.
In both domains, the system maintains epistemic distinctions that current toolchains flatten, and supports cross-institution queries that current toolchains cannot express.
\end{abstract}


% ============================================================
\section{Introduction}
\label{sec:intro}

Complex scientific and engineering work increasingly involves multiple reasoning systems operating over shared domain models.
A structural engineer designing a bracket works simultaneously with geometric constraint solvers, finite element analysis, material databases, and AI-driven generative design.
A medicinal chemist evaluating a drug candidate works with molecular docking, ADMET prediction models, experimental binding assays, and pharmacokinetic simulations.
In both cases, each tool operates with its own notion of well-formedness, its own models, and its own satisfaction criteria---yet they must share data and maintain coherent conclusions about the same domain objects.

Two developments make this heterogeneity newly urgent.
First, large language models and ML-based prediction tools are now embedded in these workflows, producing outputs that look like domain conclusions but carry no formal warranty.
An AI-predicted safety factor and a simulation-computed safety factor look identical in a spreadsheet.
An ML-predicted binding affinity and an experimentally measured IC$_{50}$ look identical in a compound database.
Second, for regulated industries---aerospace, automotive, pharmaceutical---the epistemic status of conclusions matters operationally.
Regulators need to know not just \emph{what} was concluded but \emph{how}: was it measured, predicted, simulated, or proved?

Current knowledge graph infrastructure cannot express these distinctions.
RDF triple stores~\cite{rdf} and property graphs model data as flat elements over a schema---points with no internal structure.
This is the \emph{category of elements} construction from categorical database theory: a schema is a small category $C$, an instance is a functor $F: C \to \Set$, and the total category $\int F$ contains all data as pairs (table, row).
The projection $\pi: \int F \to C$ is a discrete fibration---the fibers are sets, and elements within a fiber have no relationships to each other.

This flatness is exactly the limitation.
A finite element displacement field is not a point.
It has mesh refinement relationships (coarse-mesh and fine-mesh solutions are related by convergence), quantity extraction maps (von Mises stress, deflection, principal strain are all derived from the same field), and convergence behavior (successive refinements approach a limit).
A molecular docking pose is not a point.
It has a scoring landscape (nearby conformations have related scores), re-scoring under different force fields (the same pose evaluated differently), and ensemble relationships (multiple poses for the same ligand form a conformational ensemble).
Storing these as flat database records loses their internal mathematical life.

Goguen and Burstall's \emph{institution theory}~\cite{goguen1992} provides the repair.
An institution abstracts a logical system to signatures, sentences, models, and a satisfaction relation.
Diaconescu's \emph{Grothendieck construction}~\cite{diaconescu2002,diaconescu2008} composes a diagram of institutions into a single heterogeneous institution whose fibers are no longer sets but \emph{categories}---full logical systems with their own internal morphisms.
The move from $\int F$ (category of elements over $\Set$) to $\int \mathcal{I}$ (Grothendieck institution over $\Cat$) is precisely the generalization from flat data storage to structured, multi-logic knowledge management.

We present \textbf{Eigenius}, a platform whose architecture implements this generalization.
Its capability protocol registers domain-specific institutions that contribute structured fibers to a shared knowledge graph.
Its epistemic tracking---four categories (declared, observed, derived, verified)---arises from fiber membership in the Grothendieck construction.
We develop two worked examples spanning engineering and life sciences.


% ============================================================
\section{From Flat Fibers to Structured Fibers}
\label{sec:fibers}

\subsection{The Category of Elements (Flat Fibers)}

In categorical database theory~\cite{spivak2012}, a database schema is a small category $C$ and an instance is a functor $F: C \to \Set$.
The \emph{category of elements} $\int F$ has as objects all pairs $(c, x)$ where $c \in |C|$ and $x \in F(c)$, with morphisms induced by the schema's foreign keys.
The projection $\pi: \int F \to C$ is a \emph{discrete fibration}: the fiber over each schema node is a set (no morphisms between elements within a fiber).

This is the mathematical content of RDF, property graphs, and typed knowledge graphs like Atomic Data~\cite{atomicdata}: every datum is a point sitting over a schema node, and the entire graph is navigable by following schema-level relationships (predicates, foreign keys, property references) or by scanning fibers (all instances of a given type).

\subsection{Why Flat Fibers Are Insufficient}

When the data originates from a single uniform source (a relational database, a document store), flat fibers are appropriate.
But when data originates from multiple \emph{reasoning systems}, each with its own notion of well-formedness and its own internal structure, flatness discards essential information.

Consider three results stored in a knowledge graph as typed resources:

\begin{enumerate}
\item An FEA displacement field computed on a 5{,}000-element mesh, yielding safety factor 2.34.
\item The same analysis on a 50{,}000-element mesh, yielding safety factor 2.31.
\item The same analysis on a 500{,}000-element mesh, yielding safety factor 2.305.
\end{enumerate}

In a flat representation, these are three independent points in the fiber over \code{StressResult}.
The fact that they form a \emph{convergence sequence}---that 2.305 is a better approximation than 2.34, that the sequence is converging, that the limit is approximately 2.30---is a relationship \emph{within} the fiber that flat storage cannot express.
It is not a provenance relationship (none caused the others) and not a schema relationship (they are all the same type).
It is an internal morphism of the FEA institution.

Similarly, in drug discovery: a molecular docking study produces an ensemble of binding poses for a candidate compound.
Each pose has a docking score, but the poses are related by conformational proximity (nearby poses in RMSD space), by re-scoring relationships (the same pose evaluated under different force fields), and by clustering structure (poses grouped by binding mode).
Flattening these into independent \code{DockingResult} records loses the ensemble structure that medicinal chemists use to assess binding robustness.

\subsection{The Grothendieck Construction (Structured Fibers)}

The fix is to replace the $\Set$-valued functor with a $\Cat$-valued pseudofunctor.
Instead of each schema node having a \emph{set} of instances, it has a \emph{category} of instances---objects with morphisms between them.
The Grothendieck construction $\int \mathcal{I}$ over this pseudofunctor produces a total category whose fibers are full categories, not discrete sets.

In institution-theoretic terms~\cite{goguen1992,diaconescu2008}: an institution $\mathcal{I} = (\Sign, \Sen, \Mod, \models)$ provides, for each signature $\Sigma$, a \emph{category} of models $\Mod(\Sigma)$ (with morphisms between models) and a set of sentences $\Sen(\Sigma)$, connected by a satisfaction relation that is invariant under change of notation.
A diagram of institutions connected by comorphisms composes into a Grothendieck institution $\mathcal{I}^\#$ whose total category contains all models from all component institutions, with their internal morphisms preserved.

The fiber over the FEA institution now contains displacement fields \emph{with} their mesh refinement morphisms.
The fiber over the docking institution contains pose ensembles \emph{with} their conformational relationships.
The fiber over a formal verification institution contains proof terms \emph{with} their reduction behavior.
These are structured mathematical objects, not flat data points.


% ============================================================
\section{The Eigenius Platform}
\label{sec:eigenius}

Eigenius is a typed knowledge graph platform where all data is represented as typed resources in a uniform format called Eigon, derived from Atomic Data~\cite{atomicdata}.
We describe how its architecture implements the Grothendieck institution pattern.

\subsection{Shared Signature Category}

Eigon resources are identified by IRIs and typed by classes in an ontology organized as an immutable layer stack.
The \textbf{category of Eigon signatures} has ontology snapshots (layer stack configurations) as objects and layer extensions as morphisms.
This signature category is shared across all institutions---it is the base over which the fibers are constructed.

A single design principle governs the kernel's role: \emph{the core ontology validates; registered capabilities reason.}
The kernel enforces structural well-formedness (required properties present, types correct, namespace integrity) but has no knowledge of domain-specific logics.
Each domain logic is a registered capability---an institution that provides its own sentences, models, satisfaction relation, and (in the structured-fiber extension) its own internal morphisms.

\subsection{Eigon as Shared Signature Category}

Eigon is not an institution---it is the \textbf{shared signature category} over which the Grothendieck construction is performed.
The category of Eigon signatures has ontology snapshots as objects and layer extensions as morphisms.
The base satisfaction condition is the kernel's structural validation: 12 rules that define well-formedness of resources against class definitions.
Every institution builds over this base, adding domain-specific sentences and models to the shared Eigon foundation.

This is a mathematical necessity, not a design choice.
Making Eigon an institution would require a meta-language for institutions to register with, leading to infinite regress.
Instead, Eigon is the fixed ground---the shared language in which institutions declare their types, morphisms, and queries.

\subsection{Kernel Services vs.\ Institutions}

Two kernel services operate on the Eigon base and are \emph{not} institutions:

\textbf{Structural validation} checks resources against class definitions---required properties, type matches, format and range constraints.
This is the base satisfaction condition of the shared signature category.

\textbf{EigenTT type checking} via NbE~\cite{minitt,nbe} validates program composition---function types, dependent pairs, pattern matching.
It is necessary for the kernel to function; programs must type-check before execution.

A formal verification system such as Lean~4~\cite{lean4}, by contrast, \emph{is} an institution.
It provides the ``verified'' epistemic level by checking domain-specific mathematical propositions---proofs that a safety factor bound holds, that a binding affinity prediction satisfies a confidence interval, that a PK model is well-posed.
EigenTT checks that a program \emph{composes}; Lean~4 checks that a result is \emph{correct}.
The distinction is operational: the kernel cannot function without EigenTT; it functions without Lean~4.

\subsection{Capabilities as Institution Registrations}

Each ontology class may be associated with an evaluator, validator, and parser---collectively, a \emph{capability}.
A capability registration is an institution declaration: for a given signature, it provides the sentence functor, model functor, and satisfaction relation.

\subsection{Ground Type Resolution as Comorphism}

The kernel's ground type resolution interface---where EigenTT type checking resolves class references against the layer stack---is a kernel-internal comorphism from the Eigon structural base to the EigenTT type system.
Signatures translate forward (classes become dependent record types via NbE); models translate backward (well-typed programs reduce to valid resource instances).
This comorphism is not registered via the institution protocol---it is fixed infrastructure that all institutions depend on.

\subsection{Fiber Advertisement and Reasoning}

At registration, each institution provides a \emph{fiber declaration}: a structured description of its morphism types, query types, and advisory structural properties.
These are committed to the knowledge graph as ordinary ontology resources.
A \code{MeshRefinement} relating two \code{StressResult} resources, or a \code{ConformationalProximity} relating two \code{DockingPose} resources, is a typed resource in the knowledge graph, stored and validated by the kernel like any other resource.
The kernel knows nothing about what mesh refinement or conformational proximity \emph{means}---that is the institution's concept, governed by the institution's logic.

Each institution exposes a \emph{fiber reasoner} with four operations:

\begin{description}[style=unboxed,leftmargin=0em]
\item[\code{fiber\_declaration}] Called once at registration. Returns morphism classes, query types, and structural properties as ontology resources.
\item[\code{validate\_morphism}] Called at load time when a morphism resource enters the graph. The kernel performs structural validation (required properties, types); the institution checks domain-specific validity (e.g., that a claimed mesh refinement actually exhibits convergence).
\item[\code{query}] Called at evaluation time when a program needs fiber reasoning. The query is a typed resource (a \code{FiberQuery} subclass); the result is a typed resource.
\item[\code{discover\_morphisms}] Called on explicit request. Given resources in the institution's fiber, infer morphisms not yet in the knowledge graph. Returns candidate resources; the program decides whether to commit them.
\end{description}

Structural properties (e.g., ``\code{MeshRefinement} is a partial order'') are advisory: the kernel stores them as documentation but does not enforce them.
The institution enforces its own invariants through \code{validate\_morphism}.
This preserves the architectural boundary: the kernel validates structurally, the institution reasons within its domain.


% ============================================================
\section{Epistemic Status as Fiber Membership}
\label{sec:epistemic}

Four epistemic categories arise from the position of a resource within the Grothendieck construction:

\begin{description}[style=unboxed,leftmargin=0em]
\item[Declared] --- signature content placed by human judgment. Ontology definitions, design requirements, protocol specifications. The institution vouches for well-formedness, not truth.
\item[Observed] --- a model element imported from an external source with tracked provenance. Experimental measurements, certified material data, clinical trial outcomes.
\item[Derived] --- a model element produced by evaluation within an institution, with a reasoning trace recording the full computation. AI predictions, simulation results, extracted entities.
\item[Verified] --- a model element carrying a proof term checked by a formal institution. The proof term has internal reduction behavior (it is a structured fiber object, not a flat annotation).
\end{description}

These are not labels attached to flat terms.
They are distinguished by \emph{which fiber} of the Grothendieck construction a resource inhabits and \emph{what relationship} it bears to the institution that produced it.
Transitions are monotonic upward: derived can be promoted to verified by attaching a proof; nothing is silently downgraded.


% ============================================================
\section{Example 1: Mechanical Engineering}
\label{sec:mech}

\subsection{Domain Institutions}

\textbf{$\mathcal{I}_{\text{CAD}}$}: Signatures are parametric feature trees. Sentences are geometric constraints. Models are geometric realizations. Internal fiber morphisms: parametric variations (changing a dimension produces a related geometry), Boolean operation sequences (union/difference of features).

\textbf{$\mathcal{I}_{\text{FEA}}$}: Signatures are mesh discretizations with material and loading. Sentences are structural bounds (``von Mises $<$ yield $\times$ safety factor''). Models are displacement fields. Internal fiber morphisms: \emph{mesh refinements} (coarse $\to$ fine, with convergence delta), \emph{quantity extractions} (displacement field $\to$ scalar measurement at a point).

\textbf{$\mathcal{I}_{\text{GenAI}}$}: Signatures are design spaces. Sentences are objectives (``minimize mass subject to safety factor $\geq$ 2.0''). Models are design candidates. Internal fiber morphisms: \emph{neighborhood relationships} in parameter space, \emph{Pareto dominance} between candidates on multi-objective fronts.

\subsection{Design Session}

An engineer designs a load-bearing bracket. Design intent (5kN load, M8 bolt pattern, bounding box) enters as \textbf{declared} resources. Material properties (6061-T6 aluminum, yield 276~MPa) enter as \textbf{observed} resources from a certified database.

A generative design component produces five candidates. Candidate~3 estimates safety factor 2.4. An FEA solver computes safety factor 2.34 for the same geometry. Both are \textbf{derived}, but their provenance chains differ: GenAI passes through a non-deterministic ML component; FEA passes through a deterministic PDE solver. The system surfaces this discrepancy through a standard query over the shared knowledge graph.

For the critical bolt-hole, a Lean~4~\cite{lean4} proof establishes that the stress concentration bound holds across the material's tolerance range. This is \textbf{verified}---the proof term is a structured object in the CIC fiber with its own reduction behavior, not a flat annotation.

\subsection{Fiber Morphisms in Practice}

The FEA institution contributes three stress results connected by \code{MeshRefinement} morphisms. An EigenQL query can follow these morphisms: ``find the finest-mesh result for Candidate~3'' or ``flag results whose convergence delta exceeds 5\%.'' These queries navigate fiber structure that flat storage would discard.

The GenAI institution contributes five candidates connected by \code{ParetoDominance} morphisms on the mass-vs-strength front. A query can ask: ``which candidates are Pareto-optimal?'' This is an intra-institutional question answered by navigating the GenAI fiber, not by re-running the optimizer.


% ============================================================
\section{Example 2: Biopharmaceutical R\&D}
\label{sec:pharma}

\subsection{Domain Institutions}

\textbf{$\mathcal{I}_{\text{Dock}}$ (Molecular docking)}: Signatures are protein--ligand systems (target structure, compound library, force field parameters). Sentences are binding hypotheses (``compound X binds target Y with $\Delta G < -8$ kcal/mol''). Models are binding poses---specific spatial arrangements of ligand atoms within the binding site. Internal fiber morphisms: \emph{conformational proximity} (RMSD between poses), \emph{re-scoring} (same pose evaluated under different scoring functions), \emph{ensemble clustering} (poses grouped by binding mode).

\textbf{$\mathcal{I}_{\text{ADMET}}$ (ADMET prediction)}: Signatures are compound structures with descriptor sets. Sentences are pharmacokinetic bounds (``oral bioavailability $> 30\%$,'' ``hERG IC$_{50} > 10\,\mu$M,'' ``no CYP3A4 inhibition''). Models are predicted ADMET profiles. Internal fiber morphisms: \emph{model uncertainty} (ensemble predictions from multiple ML models for the same compound, related by agreement/disagreement), \emph{descriptor sensitivity} (how profile changes with structural modifications).

\textbf{$\mathcal{I}_{\text{Assay}}$ (Experimental assay)}: Signatures are assay protocols (target, assay format, detection method, concentration range). Sentences are activity thresholds (``IC$_{50} < 100$ nM''). Models are dose--response curves. Internal fiber morphisms: \emph{replicate relationships} (repeated measurements of the same compound under the same protocol), \emph{protocol variations} (same compound measured in different assay formats), \emph{curve fitting} (raw data points $\to$ fitted IC$_{50}$ via Hill equation).

\textbf{$\mathcal{I}_{\text{PK}}$ (Pharmacokinetic modeling)}: Signatures are compartmental models with physiological parameters. Sentences are PK targets (``$C_{\max}$ at therapeutic dose within $[1, 10]\,\mu$M,'' ``half-life $> 4$ hours''). Models are concentration--time profiles. Internal fiber morphisms: \emph{compartment refinement} (one-compartment $\to$ two-compartment model of the same data), \emph{parameter sensitivity} (how the profile changes as clearance or volume of distribution varies within confidence intervals).

\subsection{Comorphisms}

$\rho_{\text{Dock} \to \text{Assay}}$: translates a predicted binding affinity into an expected IC$_{50}$ range (signature translation), maps binding hypotheses into assay-testable predictions (sentence translation), and reduces a dose--response curve to the docking parameters it was designed to test (model reduct).

$\rho_{\text{ADMET} \to \text{PK}}$: translates predicted ADMET properties (clearance, permeability, protein binding) into compartmental model parameters. An ML-predicted clearance value becomes an input to the PK model. The comorphism tracks this dependency: when the ADMET model is retrained and predictions change, the PK institution knows which concentration--time profiles require recomputation.

\subsection{Drug Candidate Evaluation}

A medicinal chemistry team evaluates compound EIG-0042 against a kinase target.

\textbf{Step 1: Literature and databases (observed).}
The target's crystal structure (PDB entry, 1.8\,\AA\ resolution) and known SAR data from ChEMBL are loaded as observed resources. Provenance records the database, version, and retrieval date.

\textbf{Step 2: Molecular docking (derived, physics-based).}
Docking produces an ensemble of 50 poses. The top-scoring pose predicts $\Delta G = -9.2$ kcal/mol. The ensemble is stored with its conformational proximity morphisms---the medicinal chemist can query ``show me all poses within 2\,\AA\ RMSD of the top-scoring pose'' to assess binding mode robustness. The docking institution records force field, search algorithm, and exhaustiveness parameters. Epistemic status: \textbf{derived} (deterministic, physics-based scoring function).

\textbf{Step 3: ADMET prediction (derived, ML-based).}
An ML ensemble predicts oral bioavailability of 45\%, no hERG liability, but moderate CYP3A4 inhibition risk. The ensemble's individual model predictions are stored as fiber morphisms---model agreement morphisms connect the five sub-model predictions. Where sub-models disagree (CYP3A4: three predict safe, two predict risk), the disagreement is structural, not hidden. Epistemic status: \textbf{derived} (non-deterministic, ML-based).

\textbf{Step 4: Experimental validation (observed).}
A biochemical assay measures IC$_{50} = 42$ nM (three replicates: 38, 44, 45 nM). The dose--response curve enters as an observed resource. Replicate relationships are fiber morphisms within the assay institution. The measured IC$_{50}$ can be compared against the docking prediction via the comorphism $\rho_{\text{Dock} \to \text{Assay}}$: predicted $\Delta G = -9.2$ kcal/mol corresponds to an expected IC$_{50}$ of approximately 180 nM. The fourfold discrepancy between prediction (180 nM) and measurement (42 nM) is surfaced automatically.

\textbf{Step 5: PK modeling (derived, mechanistic).}
A two-compartment PK model, parameterized from ADMET predictions and experimental clearance data, predicts $C_{\max} = 3.2\,\mu$M at a 100~mg oral dose with half-life 6.8~hours. The PK institution stores the concentration--time profile with compartment refinement morphisms (a one-compartment model of the same data predicted $C_{\max} = 4.1\,\mu$M---the two-compartment refinement is a morphism in the PK fiber). Epistemic status: \textbf{derived} (deterministic, mechanistic model with ML-derived parameters).

\subsection{Cross-Institution Queries}

\emph{``Which compounds have experimentally confirmed IC$_{50}$ below 100 nM but ADMET flags from the ML ensemble?''} --- joins assay observations with ADMET predictions.

\emph{``For EIG-0042, what is the complete provenance chain from the crystal structure to the predicted therapeutic window?''} --- traverses observed (PDB structure) $\to$ derived (docking poses) $\to$ derived (ADMET profile) $\to$ derived (PK model) $\to$ derived ($C_{\max}$ prediction), crossing three institution boundaries.

\emph{``Where do the ML sub-models disagree on ADMET properties for our clinical candidates?''} --- navigates the model agreement fiber morphisms within $\mathcal{I}_{\text{ADMET}}$ to identify compounds with uncertain safety profiles, which is exactly what regulatory reviewers need to assess.


% ============================================================
\section{Implementation}
\label{sec:impl}

Five phases of Eigenius are implemented in Rust (kernel) and TypeScript (orchestration).
The Eigon data model~\cite{atomicdata} provides the shared signature category.
EigenQL (stratified Datalog with aggregation~\cite{alice}) serves as the query language.
EigenTT~\cite{minitt} dependent type checking via NbE~\cite{nbe} provides program validation as a kernel service.
IO components (LLM calls, HTTP requests) are dispatched to an orchestration layer via gRPC.
Persistent storage uses RocksDB with CBOR~\cite{rfc8949} serialization; the same key encoding supports TiKV for horizontally scalable deployments.

\subsection{NbE with Capability Modes}

A single NbE evaluator handles both type checking and program execution.
The evaluator accepts a \emph{capability mode} parameter that controls which effects are available:

\begin{description}[style=unboxed,leftmargin=0em]
\item[Pure] Standard NbE normalization. No side effects. Used by the type checker.
\item[Read] Pure plus read access to the layer chain. Used for ground type resolution.
\item[IO] Read plus component dispatch and trace production. Used for program execution.
\end{description}

In Pure mode, an IO component application produces a neutral term---it cannot reduce further without dispatch.
In IO mode, the evaluator intercepts the application, checks the trace cache for memoization, and dispatches to the orchestration layer if needed.
The result is a value paired with a reasoning trace.

This design avoids maintaining two separate evaluators (one for type checking, one for execution) while preserving the type-theoretic invariants of NbE: in Pure mode, the evaluator is total and its normal forms are canonical.
The capability mode is threaded through closures via explicit context passing, ensuring that IO dispatch occurs at the correct evaluation depth even within recursive let-bindings.

\subsection{Reasoning Traces as Proof Artifacts}

Each program execution produces a tree-structured reasoning trace that mirrors the expression tree.
After completion, the trace is pruned: only non-deterministic IO component traces are retained.
Deterministic computation (let-bindings, projections, constructions) is reconstructible from the program and the IO traces.
This follows the proofs-as-programs principle: the IO traces are axioms, the program is the theorem, and re-evaluation is the proof.

IO component traces are committed as typed resources in a \emph{trace layer}---an immutable extension of the layer chain.
The trace layer contains the program trace summary (metrics, timestamps) and all IO component traces as resources with content-addressed identifiers.
This makes reasoning traces queryable via EigenQL: ``which LLM calls contributed to this output?''

\subsection{Type Theory Extensions}

EigenTT has been extended with features inspired by Lean~4~\cite{lean4}:
propositional equality (identity type with reflexivity and J~eliminator),
decidable equality on ground types,
native constraint checking (range bounds, string patterns, format validation reduce at type-check time),
and a three-level universe stratification matching the epistemic levels.
Ground type resolution maps the full ontology---required properties as dependent pair fields, recommended properties as option types, \code{allows\_only} constraints as enumeration types, and multiple \code{class\_types} as union types.
Property access type inference walks the dependent pair chain to resolve field types.

These extensions close the gap between the Eigon ontology constraints and the EigenTT type system without introducing refinement types or subtyping.
The design follows Lean~4's approach: decidable propositions as types with native decision procedures, closed record types with explicit optional access for undeclared properties.

\subsection{Fiber Extension Path}

The current implementation uses flat fibers---resources are stored as typed Eigon values without intra-institutional morphisms.
The minimal extension to structured fibers requires no kernel changes: institutions advertise morphism types as ordinary ontology classes, and the kernel stores them like any other resource.
The richer extension---fiber reasoners as dispatched capabilities---follows the existing capability protocol pattern, with the \code{FiberReasoner} trait providing query, morphism validation, and morphism discovery methods.
Both extensions preserve the architectural boundary: the kernel validates, capabilities reason.


% ============================================================
\section{Related Work}
\label{sec:related}

\textbf{Institution theory.}
Goguen and Burstall~\cite{goguen1992} introduced institutions; Diaconescu~\cite{diaconescu2008} developed institution-independent model theory.
Mossakowski's Hets~\cite{mossakowski2007hets} is the primary implementation for heterogeneous specification.
The OMG DOL standard~\cite{ontoiop} uses institutions for ontology integration.
Meseguer's general logics~\cite{meseguer1989} and CafeOBJ~\cite{cafeobj} provide complementary formalizations.
Eigenius differs in targeting AI-driven workflows with epistemic tracking.

\textbf{Categorical databases.}
Spivak's functorial data model~\cite{spivak2012} formalizes databases as $\Set$-valued functors.
The category of elements construction is the flat-fiber case that Eigenius generalizes.

\textbf{Knowledge graphs.}
RDF~\cite{rdf} and property graphs are operationally categories of elements with flat fibers.
Eigenius's Eigon format derives from Atomic Data~\cite{atomicdata}, adapted with non-fetchable IRIs, multiple class membership, and a three-layer type system.

\textbf{AI in engineering and drug discovery.}
Generative design tools and ML-based ADMET predictors produce results with no formal epistemic tracking.
Eigenius makes the distinction between AI prediction, simulation, experimental observation, and formal proof a structural property.

\textbf{Neurosymbolic architectures.}
DeepProbLog~\cite{deepproblog} and Logic Tensor Networks~\cite{ltn} integrate neural and symbolic computation at the gradient level.
Eigenius uses institution theory rather than differentiable logic as the integration framework, treating AI components as black boxes in typed pipelines.


% ============================================================
\section{Discussion}
\label{sec:discussion}

\subsection{The Satisfaction Condition as Architectural Invariant}

The institution-theoretic satisfaction condition---truth is invariant under change of notation---maps to a key Eigenius invariant: validation results are invariant under layer extension.
Adding ontology definitions cannot invalidate existing validated resources.
For domain institutions, the satisfaction condition must be maintained by the capability implementation---a testable property via property-based testing.

\subsection{When Flat Fibers Suffice}

Not every query requires structured fibers.
``What is the safety factor for Candidate~3?'' and ``what is the IC$_{50}$ for EIG-0042?'' are answerable from flat storage with provenance traces.
Structured fibers become necessary when queries navigate \emph{within} an institution: convergence analysis, ensemble robustness, model agreement, parameter sensitivity.
The architecture supports both: flat fibers are the discrete special case of structured fibers.

\subsection{Why Not a Join Table?}

A natural objection: fiber morphisms could be represented as relational links---a join table connecting pairs of resources with a relationship type.
This is what RDF triples and property graph edges provide.
The Grothendieck construction adds four things that relational links cannot express.

First, \emph{morphism validation requires domain reasoning}.
A join table records that two stress results are related by ``mesh refinement.''
But whether the relationship is \emph{valid}---whether the fine mesh actually refines the coarse mesh, whether boundary conditions match, whether the convergence delta is positive---is a domain-specific judgment that the FEA institution makes through \code{validate\_morphism}.
A relational link carries no satisfaction condition.

Second, \emph{morphisms compose with domain semantics}.
Three mesh refinements in sequence form a convergence chain with a computable convergence \emph{rate}.
A join table records three pairs; the institution knows they compose and can compute the rate, detect convergence, or flag divergence.
Composition is a categorical operation, not a relational one.

Third, \emph{fiber queries navigate with domain knowledge}.
``Find the finest-mesh result whose convergence delta is below 5\%'' requires interpreting convergence within the FEA fiber.
EigenQL can match the resources structurally, but the institution provides the semantic evaluation.

Fourth, \emph{morphism discovery is an institution capability}.
Given five stress results, the FEA institution can \emph{infer} which ones are refinements of which, based on mesh topology and boundary conditions, without anyone manually populating the links.
A join table requires external population; a fiber reasoner generates its own morphisms.

In short: a relational link records \emph{that} two results are related.
A fiber morphism captures \emph{how} they are related, \emph{whether} the relationship is valid, and \emph{what} can be derived from it.
The morphism is not a link---it is a typed, validated, queryable object with domain semantics, governed by the institution that produced it.

\subsection{Epistemic Status Across the Dispatch Boundary}

If institutions are black boxes behind a dispatch interface, how does the kernel determine the epistemic status of their results?
Three mechanisms work together, none requiring the kernel to trust the institution.

First, the \emph{trace provides the default}.
Every capability dispatch produces a reasoning trace recording the component, inputs, and output.
Any result produced by dispatch is at least \textbf{derived}---the kernel knows a computation occurred because it initiated the dispatch.
This requires no institution cooperation.

Second, the \emph{institution may assert stronger status} by including epistemic base classes in the returned resources.
A formal verification institution may return a resource with \code{is\_a} including \code{VerifiedResource} and a \code{verification} property pointing to a checked proof term.
The kernel does not take this claim on trust---it structurally validates that the required provenance properties are present (\code{VerifiedResource} requires \code{verification}; \code{ObservedResource} requires \code{source}).
An institution that claims ``verified'' without providing a proof term fails structural validation.

Third, the \emph{provenance chain records input heterogeneity}.
A fiber reasoner computing a convergence report may consume inputs at different epistemic levels---an FEA result (derived, deterministic), a GenAI estimate (derived, non-deterministic), and a material property (observed).
The trace records which resources were consulted; each already carries its own epistemic status.
A downstream query can traverse the provenance chain to reveal this heterogeneity without the fiber reasoner encoding it explicitly.

The result is that epistemic status flows through the dispatch boundary without special machinery: traces provide the floor, base classes allow upgrades subject to structural validation, and provenance chains preserve the full picture.

\subsection{Limitations}

AI-based institutions ($\mathcal{I}_{\text{GenAI}}$, $\mathcal{I}_{\text{ADMET}}$) are degenerate in that their satisfaction relations are not decidable.
Institution theory accommodates this---decidability is not required---but their guarantees are weaker.
The epistemic system makes this weakness \emph{visible}: ML-derived results carry different provenance than simulation or experimental results.

The formalization is incomplete: capability implementations are trusted to satisfy the institution axioms.
Property-based testing provides confidence but not proof.

% ============================================================
\section{Conclusion}
\label{sec:conclusion}

The move from flat data storage (the category of elements over $\Set$) to structured multi-logic knowledge management (the Grothendieck institution over $\Cat$) is the mathematical content of what Eigenius implements operationally.
Each domain institution contributes a fiber with internal structure---mesh convergence, conformational ensembles, proof reductions, model agreement---that flat knowledge graphs discard.
Epistemic categories arise from fiber membership: which institution produced a resource, and what kind of object it is within that institution's fiber.

The two worked examples demonstrate that the pattern is domain-agnostic: the same architecture handles mechanical engineering (CAD, FEA, generative design) and biopharmaceutical R\&D (docking, ADMET, assays, PK modeling).
In both cases, the system maintains epistemic distinctions that current toolchains flatten, supports cross-institution queries that current toolchains cannot express, and preserves the internal structure of domain results that flat storage loses.

The guiding architectural principle---the core ontology validates, registered capabilities reason---is an institution-theoretic statement: the kernel provides the shared signature category and enforces satisfaction structurally; each capability provides its own logic.
The Grothendieck construction is what this principle produces when taken to its mathematical conclusion.

\section*{Acknowledgments}

The author thanks Allen Brown for calling out the importance and possibility of structured fibrations beyond the flat $\Set$ category---an insight he attributes to a conversation with David Spivak~\cite{spivak2012}---that results from domain-specific reasoning systems carry internal mathematical structure that $\Set$-valued storage discards, and that the Grothendieck construction over $\Cat$ provides the precise generalization needed to preserve it.

% ============================================================
\appendix
\section{Fiber Reasoner Protocol Sketch}
\label{app:fiber}

This appendix presents the \code{FiberReasoner} trait---the interface by which an institution exposes reasoning about its internal fiber structure to the kernel.
The design follows the existing capability protocol: the kernel dispatches, the institution reasons.
The full specification is in design document D10 (Grothendieck Institution Protocol).

\subsection{Trait Definition}

{\scriptsize
\begin{verbatim}
trait FiberReasoner: Send + Sync {
  /// Execute a fiber query (typed Eigon resource,
  /// subclass of FiberQuery). Institution defines
  /// its own query types.
  fn query(
    &self, query: &Resource, ctx: &ExecutionContext,
  ) -> Result<Resource, Error>;

  /// Validate a claimed morphism against the
  /// institution's domain logic (not structural
  /// validation — the kernel handles that).
  fn validate_morphism(
    &self, morphism: &Resource,
    ctx: &ExecutionContext,
  ) -> Result<MorphismValidation, Error>;

  /// Given resources in this institution's fiber,
  /// infer morphisms not yet in the knowledge graph.
  /// Returns resources; caller decides storage.
  fn discover_morphisms(
    &self, resources: &[Resource],
    ctx: &ExecutionContext,
  ) -> Result<Vec<Resource>, Error>;

  /// Declare morphism types, query types, and
  /// structural properties. Called once during
  /// capability registration.
  fn fiber_declaration(
    &self,
  ) -> FiberDeclaration;
}
\end{verbatim}
}

Each method has a distinct role in the system lifecycle:
\code{fiber\_declaration} is called at registration (the institution advertises its fiber);
\code{validate\_morphism} at load time (a morphism enters the graph);
\code{query} at evaluation time (a program needs fiber reasoning);
\code{discover\_morphisms} on explicit request by a program.

\subsection{Fiber Declaration}

During registration, the institution declares its fiber structure as ordinary ontology resources:

{\scriptsize
\begin{verbatim}
struct FiberDeclaration {
  /// The institution's IRI
  institution_iri: Iri,
  /// Human-readable name
  name: String,

  /// Morphism classes (e.g., MeshRefinement,
  /// ConformationalProximity), each with
  /// source type, target type, properties.
  morphism_types: Vec<Resource>,

  /// FiberQuery subclasses this institution
  /// can answer (e.g., ConvergenceQuery).
  query_types: Vec<Resource>,

  /// Advisory structural properties (e.g.,
  /// "MeshRefinement is a partial order").
  /// Kernel stores but does not enforce.
  structural_properties: Vec<Resource>,
}
\end{verbatim}
}

Structural properties are advisory: declaring that \code{MeshRefinement} is transitive does not cause the kernel to enforce transitivity.
The institution enforces its own properties through \code{validate\_morphism}.
The advisory declarations serve as documentation and as future optimization hints (if the kernel knows a morphism type is transitive, it can answer reachability queries without dispatching to the fiber reasoner for intermediate steps).

\subsection{Dispatch Model}

The kernel's role is minimal:
\begin{enumerate}
\item At registration: call \code{fiber\_declaration()}, store the returned types as ontology resources.
\item At load time: structurally validate morphism resources (kernel's job), optionally dispatch to \code{validate\_morphism()} for domain validation (institution's job).
\item At evaluation time: recognize \code{FiberQuery} subclasses, dispatch to \code{query()}, return typed results.
\item \code{discover\_morphisms()} is called by programs, not by the kernel automatically. The fiber reasoner has no write access---it returns resources, and the program decides whether to commit them.
\end{enumerate}

The kernel never interprets morphism semantics.
It never computes convergence, RMSD, Pareto dominance, or definitional equality.
It stores, validates structurally, dispatches, and returns results.

\subsection{EigenQL Integration}

No query language changes are needed.
Fiber morphisms are ordinary resources, so existing EigenQL queries work:

{\scriptsize
\begin{verbatim}
USING "urn:eigenius:fea:MeshRefinement",
      "urn:eigenius:fea:StressResult"
MATCH MeshRefinement(?m) {
  source: ?s1, target: ?s2,
  convergence_delta: ?delta
},
StressResult(?s2) { safety_factor: ?sf }
WHERE ?delta > 0.05
RETURN [] {
  result: ?s2, factor: ?sf, delta: ?delta
}
\end{verbatim}
}

\noindent This finds stress results whose mesh convergence delta exceeds 5\%---``results that are not yet mesh-converged.''
For queries requiring domain reasoning, the evaluator recognizes \code{FiberQuery} subclasses and dispatches to the fiber reasoner, following the standard capability dispatch pattern.

\subsection{WASM Compatibility}

The trait maps directly to a WASM export interface: each method becomes an exported function with CBOR-encoded input and output, subject to the same memory isolation, fuel-bounded execution, and controlled interface surface as other WASM capabilities (Section~\ref{sec:eigenius}).

\balance

% ============================================================
\bibliographystyle{plainnat}
\begin{thebibliography}{18}

\bibitem[Goguen and Burstall(1992)]{goguen1992}
J.A.~Goguen and R.M.~Burstall.
\newblock Institutions: Abstract model theory for specification and programming.
\newblock \emph{Journal of the ACM}, 39(1):95--146, 1992.

\bibitem[Diaconescu(2002)]{diaconescu2002}
R.~Diaconescu.
\newblock Grothendieck institutions.
\newblock \emph{Applied Categorical Structures}, 10(4):383--402, 2002.

\bibitem[Diaconescu(2008)]{diaconescu2008}
R.~Diaconescu.
\newblock \emph{Institution-independent Model Theory}.
\newblock Birkh\"{a}user, Basel, 2008.

\bibitem[Mossakowski et~al.(2007)]{mossakowski2007hets}
T.~Mossakowski, C.~Maeder, and K.~L\"{u}ttich.
\newblock The heterogeneous tool set, {Hets}.
\newblock In \emph{TACAS 2007}, LNCS 4424, pages 519--522. Springer, 2007.

\bibitem[Meseguer(1989)]{meseguer1989}
J.~Meseguer.
\newblock General logics.
\newblock In \emph{Logic Colloquium~'87}, volume 129, pages 275--329. Elsevier, 1989.

\bibitem[Diaconescu and Futatsugi(1998)]{cafeobj}
R.~Diaconescu and K.~Futatsugi.
\newblock \emph{CafeOBJ Report}.
\newblock AMAST Series in Computing, vol.~6. World Scientific, 1998.

\bibitem[OMG(2018)]{ontoiop}
Object Management Group.
\newblock Distributed Ontology, Modeling and Specification Language ({DOL}), Version 1.0.
\newblock OMG Document formal/2018-08-01, 2018.

\bibitem[Spivak(2012)]{spivak2012}
D.I.~Spivak.
\newblock Functorial data migration.
\newblock \emph{Information and Computation}, 217:31--51, 2012.

\bibitem[Cyganiak et~al.(2014)]{rdf}
R.~Cyganiak, D.~Wood, and M.~Lanthaler.
\newblock {RDF 1.1} concepts and abstract syntax.
\newblock W3C Recommendation, 2014.

\bibitem[Meindertsma(2020)]{atomicdata}
J.~Meindertsma.
\newblock {Atomic Data}.
\newblock W3C Community Group Specification, 2020--2026.
\newblock \url{https://docs.atomicdata.dev/}.

\bibitem[Coquand et~al.(2009)]{minitt}
T.~Coquand, Y.~Kinoshita, B.~Nordstr\"{o}m, and M.~Takeyama.
\newblock A simple type-theoretic language: {EigenTT}.
\newblock In \emph{From Semantics to Computer Science}, pages 139--164. Cambridge University Press, 2009.

\bibitem[Abel et~al.(2007)]{nbe}
A.~Abel, T.~Coquand, and P.~Dybjer.
\newblock Normalization by evaluation for {Martin-L\"{o}f} type theory with typed equality judgements.
\newblock In \emph{LICS 2007}, pages 3--12. IEEE, 2007.

\bibitem[de~Moura and Ullrich(2021)]{lean4}
L.~de~Moura and S.~Ullrich.
\newblock The {Lean~4} theorem prover and programming language.
\newblock In \emph{CADE-28}, LNCS 12699, pages 625--635. Springer, 2021.

\bibitem[Abiteboul et~al.(1995)]{alice}
S.~Abiteboul, R.~Hull, and V.~Vianu.
\newblock \emph{Foundations of Databases}.
\newblock Addison-Wesley, 1995.

\bibitem[Bormann and Hoffman(2020)]{rfc8949}
C.~Bormann and P.~Hoffman.
\newblock Concise binary object representation ({CBOR}).
\newblock RFC 8949, IETF, 2020.

\bibitem[Manhaeve et~al.(2018)]{deepproblog}
R.~Manhaeve, S.~Dumancic, A.~Kimmig, T.~Demeester, and L.~De~Raedt.
\newblock {DeepProbLog}: Neural probabilistic logic programming.
\newblock In \emph{NeurIPS 2018}, 2018.

\bibitem[Badreddine et~al.(2022)]{ltn}
S.~Badreddine, A.~Garcez, L.~Serafini, and M.~Spranger.
\newblock Logic tensor networks.
\newblock \emph{Artificial Intelligence}, 303, 2022.

\end{thebibliography}

\end{document}